\begin{frame}{Modo \textit{Sample} ou \textit{Track}}
  \begin{columns}[T]
    \begin{column}{0.52\textwidth}
      \begin{block}{Chave fechada}
        A entrada fica conectada ao capacitor. O capacitor carrega ou descarrega até acompanhar a tensão de entrada.
      \end{block}
      \[ V_{out} \approx V_{in} \]
      Também é chamado de modo \textit{track}, pois a saída ``segue'' a entrada.
    \end{column}
    \begin{column}{0.44\textwidth}
      \resizebox{\linewidth}{!}{%
      \begin{tikzpicture}[scale=0.72,>=Latex]
        \draw[->] (0,0) -- (4.7,0) node[right] {tempo};
        \draw[->] (0,0) -- (0,3.3) node[above] {tensão};
        \draw[blue, thick, samples=50, domain=0:4.4] plot (\x,{1.6+0.7*sin(160*\x)});
        \draw[red, dashed, thick, samples=50, domain=0:4.4] plot (\x,{1.6+0.7*sin(160*\x)});
        \node[blue] at (3.8,2.9) {$V_{in}$};
        \node[red] at (3.8,0.55) {$V_{out}$};
      \end{tikzpicture}%
      }
    \end{column}
  \end{columns}
\end{frame}
